% Figura acumulativa; destino figures/monodromia_regional_ampliacion.tex.
\begin{figure}[p]
\centering
\begin{tikzpicture}[x=1cm,y=1cm,>=Latex,font=\small,
  flow/.style={->,line width=.55pt},
  ann/.style={align=center,font=\scriptsize}]
\node[font=\bfseries] at (6.9,9.2)
  {Conexión nonádica y coordenadas regionales};
\begin{scope}[shift={(3.4,6.8)}]
 \foreach \k in {1,...,9}{
  \pgfmathsetmacro{\ang}{90-40*(\k-1)}
  \coordinate (p\k) at (\ang:1.52);
  \fill (p\k) circle (1.4pt);
  \node[fill=white,inner sep=1.1pt,font=\scriptsize]
    at (\ang:1.83) {\(\k\)};
 }
 \foreach \k in {1,...,8}{
  \pgfmathtruncatemacro{\next}{\k+1}
  \draw[flow,shorten >=3pt,shorten <=3pt] (p\k)--(p\next);
 }
 \draw[flow,shorten >=3pt,shorten <=3pt] (p9)--(p1);
 \node[ann] at (0,.1) {\(g_k\in C_9\)\\\(\Gamma_{9,k}\)};
 \node[ann] at (0,-2.23) {Una vuelta: \(g\mapsto g\), \(m\mapsto m+1\)};
\end{scope}
\node[ann,text width=5.25cm] at (10.1,7.75)
 {\(B_k=(u_k^\pi,u_k^e,u_k^\varphi)\)\\
 bloque regional en una fibra};
\node at (8.55,6.5) {\(u^\pi\)};
\node at (11.65,6.5) {\(u^e\)};
\draw[<->,line width=.65pt] (8.95,6.5)--node[above]{\(\mathfrak c\)}(11.25,6.5);
\draw[densely dashed,gray] (10.1,5.72)--(10.1,7.23);
\node[fill=white,inner sep=2pt] at (10.1,5.6){\(u^\varphi\)};
\node[ann] at (10.1,4.6)
 {Eje \(u^\varphi\) y agregado \(u^\pi+u^e\)\\
 invariantes bajo \(\mathfrak c\)};
\draw[flow] (5.45,6.8)--(7.25,6.8);
\draw[gray!60,line width=.3pt] (.55,4.02)--(13.25,4.02);
\node[ann,text width=12cm] (hist) at (6.9,3.47)
 {Historia compatible: prefijos, cilindros, orientación, cocientes,
 marcos de incidencia y memoria};
\coordinate (fork) at (6.9,2.85);
\draw[flow] (hist.south)--(fork);
\draw[flow] (fork)--(3.6,2.2);
\draw[flow] (fork)--(10.3,2.2);
\node[ann,text width=5.6cm] at (3.6,1.65)
 {Realización arquimediana\\
 \(I_{n+1}\subseteq I_n,\quad |I_n|=729^{-n}\)\\
 límite de los prefijos regionales};
\node[ann,text width=5.6cm] at (10.3,1.65)
 {Realización incidencial\\
 \(r_{n,m}Q_n=Q_mp_{n,m}\)\\
 palabras codificadas y marcos compatibles};
\node[ann,text width=5.6cm] at (3.6,.42)
 {\(\pi,\varphi,e\); determinación dodecafásica de \(\alpha\)};
\node[ann,text width=5.6cm] at (10.3,.42)
 {Witt--Golay--Leech--\(V^\natural\)\\
 automorfismos de la realización excepcional};
\draw[flow] (3.6,1.03)--(3.6,.78);
\draw[flow] (10.3,1.03)--(10.3,.78);
\end{tikzpicture}
\caption{Nueve fases de una conexión y dos realizaciones de la historia
compatible. El ciclo representa la fase; cada vuelta incrementa el contador.
El diagrama regional representa la involución \(\mathfrak c\), que
intercambia clausura y propagación y fija el eje de autoescala en una fibra.
Esta invariancia no impone que el eje sea constante durante el transporte.
Las dos aplicaciones inferiores utilizan el mismo prefijo: una determina
su límite numérico y la otra conserva palabras y marcos de incidencia. Los
mapas de esta segunda realización se especifican en la sección
\ref{exc:doble-lectura}; el recorrido excepcional se desarrolla después.}
\label{mono-ampl:fig-regional}
\end{figure}
